% !TeX spellcheck = en_US
%!TEX root=../Main.tex
% \clearpage
\section{Concluding Remarks}

In this paper, we propose a versatile and training-efficient video foundation model \model. To our best knowledge, \model~is the first work to perform best among existing researches on all action understanding, video-language alignment, and video open understanding tasks. Compared with previous related work \cite{maskfeat,mtv,merlot}, it greatly lifts the generality of video foundation models to a new level, by achieving state-of-the-art performance on nearly 40 datasets covering 10 different tasks. 
The model exploits a unified video representation based on the cross-model learning between masked video learning (VideoMAE) and video-language contrastive modeling along with supervised training. Compared with previous foundation models, it is efficient in training. With simple ViT and its corresponding variants, we achieve generalized video representations with 64.5K GPU hours (A100-80G), while CoCa \cite{coca} requires 245.76K TPU hours (v4). We validate such generalized spatiotemporal representation on a spectrum of applications. With simple task heads (even linear ones) and proper downstream adaption tuning, our video representation demonstrates record-breaking results in all used datasets. Even for zero-shot and open-set settings, our model spectrum still gives consistent and non-trivial performance increases, further proving its generalization and adaption.

\subsection{Limitations} 
Our study shows the effectiveness and feasibility of video foundation models instead of giving brand-new formulations or model designs. It focuses on the current popular video perception tasks and handles videos using clips. Its devise can hardly process long-term video tasks, as well as high-order ones, \eg~anticipating plots from the seen parts of a movie. Gaining the capacity to address these tasks is crucial to further push the generality of video representation learning.


\subsection{Future Work}
To further extend the generality of the video foundation models, we suppose embracing model coordination and cognition is necessary for its studies. Specifically, how to systematically coordinate foundation models trained from different modalities, pretraining tasks, and even varied architectures for a better representation remains open and challenging. There are multiple technical routes to address it, \eg\ model distillation, unifying different pretraining objectives, and feature alignment, to name a few. By exploiting previously learned knowledge, we can accelerate video foundation model development sustainably. 

In the long run, foundation models are expected to master cognitive capabilities more than perceivable ones. Considering its feasibility, we suppose one of its research trends is to achieve large-scale spatiotemporal analysis (long-term \& big scene) from the foundational dynamic perception in the open world, leading to essential cognitive understanding.
Moreover, it has raised a tide that combining foundation models with decision-making to form intelligent agents to explore new tasks. In this interaction, data collection and model training are also automated. The whole process enters a closed loop as the interactive results will adjust agent strategies and behaviors. Our initial experiments (Section \ref{sec_vlt}) on vision-language navigation demonstrate the promising future of integrating video foundation models into Embodied AI.
